\documentclass[12pt, titlepage]{article}

\usepackage{booktabs}
\usepackage{tabularx}
\usepackage{hyperref}
\usepackage{longtable}
\usepackage{pdflscape}
\hypersetup{
    colorlinks,
    citecolor=blue,
    filecolor=black,
    linkcolor=red,
    urlcolor=blue
}
% \usepackage[numbers]{natbib}

\input{../Comments.text}
\input{../Common.text}


%System Test Macro Definition

\newcounter{SFT}[subsubsection]
\newcounter{SNFT}[subsubsection]
\renewcommand{\theSFT}{SFT-\arabic{subsubsection}.\arabic{SFT}}
\renewcommand{\theSNFT}{SNFT-\arabic{subsubsection}.\arabic{SNFT}}
\renewcommand{\theHSFT}{\thesection.\thesubsection.\thesubsubsection.\arabic{SFT}}
\renewcommand{\theHSNFT}{\thesection.\thesubsection.\thesubsubsection.\arabic{SNFT}}

% One labelled field (bold heading)
\newcommand{\testfield}[2]{\par\noindent\textbf{#1:} #2\par\vspace{4pt}}

% Functional test:
% {label}{name}{control}{initial state}{input}{output}{derivation}{how 
%performed}
\newcommand{\functest}[8]{%
	\par\noindent\refstepcounter{SFT}\label{#1}%
	\textbf{\theSFT} #2\par\vspace{8pt}
	\testfield{Control}{#3}
	\testfield{Initial State}{#4}
	\testfield{Input}{#5}
	\testfield{Output}{#6}
	\testfield{Test Case Derivation}{#7}
	\testfield{How test will be performed}{#8}
}

% Nonfunctional test:
% {label}{name}{type}{initial state}{input/condition}{output/result}{how 
%performed}
\newcommand{\nonfunctest}[7]{%
	\par\noindent\refstepcounter{SNFT}\label{#1}%
	\textbf{\theSNFT} #2\par\vspace{8pt}
	\testfield{Type}{#3}
	\testfield{Initial State}{#4}
	\testfield{Input/Condition}{#5}
	\testfield{Output/Result}{#6}
	\testfield{How test will be performed}{#7}
}

% --- HOW TO USE IT ---

% Functional test (inside a \subsubsection under "Tests for Functional 
%Requirements"):
%\item \functest{label}
%{Name}
%{Control: Manual or Automatic}
%{Initial State}
%{Input}
%{Output}
%{Test Case Derivation}
%{How test will be performed}

% Nonfunctional test (inside a \subsubsection under "Tests for Nonfunctional 
%Requirements"):
%\item \nonfunctest{label}
%{Name}
%{Type: Functional, Dynamic, Manual, Static, etc.}
%{Initial State}
%{Input/Condition}
%{Output/Result}
%{How test will be performed}

% TRACE TABLE MACRO 
\newcommand{\traceentry}[3]{
  #1 & #2 & #3 \\
  \hline
}


\begin{document}

\title{System Verification and Validation Plan for \progname{}} 
\author{\authname}
\date{\today}
	
\maketitle

\pagenumbering{roman}

\section*{Revision History}

\begin{tabularx}{\textwidth}{p{3cm}p{2cm}X}
\toprule {\bf Date} & {\bf Version} & {\bf Notes}\\
\midrule
Date 1 & 1.0 & Notes\\
Date 2 & 1.1 & Notes\\
\bottomrule
\end{tabularx}

~\\
\wss{The intention of the VnV plan is to increase confidence in the software.
However, this does not mean listing every verification and validation technique
that has ever been devised.  The VnV plan should also be a \textbf{feasible}
plan. Execution of the plan should be possible with the time and team available.
If the full plan cannot be completed during the time available, it can either be
modified to ``fake it'', or a better solution is to add a section describing
what work has been completed and what work is still planned for the future.}

\wss{The VnV plan is typically started after the requirements stage, but before
the design stage.  This means that the sections related to unit testing cannot
initially be completed.  The sections will be filled in after the design stage
is complete.  the final version of the VnV plan should have all sections filled
in.}

\newpage

\tableofcontents

\listoftables

\listoffigures
\wss{Remove this section if it isn't needed}

\newpage

\section{Symbols, Abbreviations, and Acronyms}

\renewcommand{\arraystretch}{1.2}
\begin{tabular}{l l} 
  \toprule		
  \textbf{symbol} & \textbf{description}\\
  \midrule 
  T & Test\\
  \bottomrule
\end{tabular}\\

\wss{symbols, abbreviations, or acronyms --- you can simply reference the SRS
  \cite{SRS} tables, if appropriate}

\wss{Remove this section if it isn't needed}

\newpage

\pagenumbering{arabic}

This document ... \wss{provide an introductory blurb and roadmap of the
  Verification and Validation plan}

\section{General Information}

\subsection{Summary}

\wss{Say what software is being tested.  Give its name and a brief overview of
  its general functions.}

\subsection{Objectives}

\wss{State what is intended to be accomplished.  The objective will be around
  the qualities that are most important for your project.  You might have
  something like: ``build confidence in the software correctness,''
  ``demonstrate adequate usability.'' etc.  You won't list all of the qualities,
  just those that are most important.}

\wss{You should also list the objectives that are out of scope.  You don't have 
the resources to do everything, so what will you be leaving out.  For instance, 
if you are not going to verify the quality of usability, state this.  It is also 
worthwhile to justify why the objectives are left out.}

\wss{The objectives are important because they highlight that you are aware of 
limitations in your resources for verification and validation.  You can't do everything, 
so what are you going to prioritize?  As an example, if your system depends on an 
external library, you can explicitly state that you will assume that external library 
has already been verified by its implementation team.}

\subsection{Custom Documents}

\wss{Summarize the CDs (Custom Documents) that were tackled by this project.  CDs 
can include usability testing, code walkthroughs, user documentation, formal
proof, GenderMag personas, Design Thinking, etc.  CDs should have already
been approved by the course instructor as included in your problem statement.
You can use a pull request to update your CDs (in TeamComposition.csv or
Repos.csv) if your plan changes as a result of the VnV planning exercise.}

\subsection{Relevant Documentation}

\wss{Reference relevant documentation.  This will definitely include your SRS
  and your other project documents (design documents, like MG, MIS, etc).  You
  can include these even before they are written, since by the time the project
  is done, they will be written.  You can create BibTeX entries for your
  documents and within those entries include a hyperlink to the documents.}

Citation Example of SRS\cite{SRS}\\

\wss{Don't just list the other documents.  You should explain why they are relevant and 
how they relate to your VnV efforts.}

\section{Plan}

\wss{Introduce this section.  You can provide a roadmap of the sections to
  come.}

\subsection{Verification and Validation Team}

\wss{Your teammates.  Maybe your supervisor.
  You should do more than list names.  You should say what each person's role is
  for the project's verification.  A table is a good way to summarize this information.}

\subsection{SRS Verification}

\wss{List any approaches you intend to use for SRS verification.  This may
  include ad hoc feedback from reviewers, like your classmates (like your
  primary reviewer), or you may plan for something more rigorous/systematic.}

\wss{If you have a supervisor for the project, you shouldn't just say they will
read over the SRS.  You should explain your structured approach to the review.
Will you have a meeting?  What will you present?  What questions will you ask?
Will you give them instructions for a task-based inspection?  Will you use your
issue tracker?}

\wss{Maybe create an SRS checklist?}

\subsection{Design Verification}

\wss{Plans for design verification}

\wss{The review will include reviews by your classmates}

\wss{Create a checklists?}

\subsection{Verification and Validation Plan Verification}

\wss{The verification and validation plan is an artifact that should also be
verified.  Techniques for this include review and mutation testing.}

\wss{The review will include reviews by your classmates}

\wss{Create a checklists?}

\subsection{Implementation Verification}

\wss{You should at least point to the tests listed in this document and the unit
  testing plan.}

\wss{In this section you would also give any details of any plans for static
  verification of the implementation.  Potential techniques include code
  walkthroughs, code inspection, static analyzers, etc.}

\wss{The final class presentation in CAS 741 could be used as a code
walkthrough.  There is also a possibility of using the final presentation (in
CAS741) for a partial usability survey.}

\subsection{Automated Testing and Verification Tools}

\wss{What tools are you using for automated testing.  Likely a unit testing
  framework and maybe a profiling tool, like ValGrind.  Other possible tools
  include a static analyzer, make, continuous integration tools, test coverage
  tools, etc.  Explain your plans for summarizing code coverage metrics.
  Linters are another important class of tools.  For the programming language
  you select, you should look at the available linters.  There may also be tools
  that verify that coding standards have been respected, like flake9 for
  Python.}

\wss{If you have already done this in the development plan, you can point to
that document.}

\wss{The details of this section will likely evolve as you get closer to the
  implementation.}

\subsection{Software Validation}

\wss{If there is any external data that can be used for validation, you should
  point to it here.  If there are no plans for validation, you should state that
  here.}

\wss{You might want to use review sessions with the stakeholder to check that
the requirements document captures the right requirements.  Maybe task based
inspection?}

\wss{For those capstone teams with an external supervisor, the Rev 0 demo should 
be used as an opportunity to validate the requirements.  You should plan on 
demonstrating your project to your supervisor shortly after the scheduled Rev 0 demo.  
The feedback from your supervisor will be very useful for improving your project.}

\wss{For teams without an external supervisor, user testing can serve the same purpose 
as a Rev 0 demo for the supervisor.}

\wss{This section might reference back to the SRS verification section.}

\section{System Tests}

\wss{There should be text between all headings, even if it is just a roadmap of
the contents of the subsections.}

\subsection{Tests for Functional Requirements}

\wss{Subsets of the tests may be in related, so this section is divided into
  different areas.  If there are no identifiable subsets for the tests, this
  level of document structure can be removed.}

\wss{Include a blurb here to explain why the subsections below
  cover the requirements.  References to the SRS would be good here.}

\subsubsection{Area of Testing1}

\wss{It would be nice to have a blurb here to explain why the subsections below
  cover the requirements.  References to the SRS would be good here.  If a section
  covers tests for input constraints, you should reference the data constraints
  table in the SRS.}
		
\paragraph{Title for Test}

\begin{enumerate}

\item{test-id1\\}

Control: Manual versus Automatic
					
Initial State: 
					
Input: 
					
Output: \wss{The expected result for the given inputs.  Output is not how you
are going to return the results of the test.  The output is the expected
result.}

Test Case Derivation: \wss{Justify the expected value given in the Output field}
					
How test will be performed: 
					
\item{test-id2\\}

Control: Manual versus Automatic
					
Initial State: 
					
Input: 
					
Output: \wss{The expected result for the given inputs}

Test Case Derivation: \wss{Justify the expected value given in the Output field}

How test will be performed: 

\end{enumerate}

\subsubsection{User Interface Tests for Functional Requirements}

\input{UISFT.text}

\subsection{Tests for Nonfunctional Requirements}

\wss{The nonfunctional requirements for accuracy will likely just reference the
  appropriate functional tests from above.  The test cases should mention
  reporting the relative error for these tests.  Not all projects will
  necessarily have nonfunctional requirements related to accuracy.}

\wss{For some nonfunctional tests, you won't be setting a target threshold for
passing the test, but rather describing the experiment you will do to measure
the quality for different inputs.  For instance, you could measure speed versus
the problem size.  The output of the test isn't pass/fail, but rather a summary
table or graph.}

\wss{Tests related to usability could include conducting a usability test and
  survey.  The survey will be in the Appendix.}

\wss{Static tests, review, inspections, and walkthroughs, will not follow the
format for the tests given below.}

\wss{If you introduce static tests in your plan, you need to provide details.
How will they be done?  In cases like code (or document) walkthroughs, who will
be involved? Be specific.}

\subsubsection{Area of Testing1}
		
\paragraph{Title for Test}

\begin{enumerate}

\item{test-id1\\}

Type: Functional, Dynamic, Manual, Static etc.
					
Initial State: 
					
Input/Condition: 
					
Output/Result: 
					
How test will be performed: 
					
\item{test-id2\\}

Type: Functional, Dynamic, Manual, Static etc.
					
Initial State: 
					
Input: 
					
Output: 
					
How test will be performed: 

\end{enumerate}

\subsubsection{User Interface Tests for Non-Functional Requirements}

\input{UISNFT.text}

\subsection{Traceability Between Test Cases and Requirements}

% \wss{Provide a table that shows which test cases are supporting which
%   requirements.}
\input{Section4TraceTable.text}

\section{Unit Test Description}

TBD

% \wss{This section should not be filled in until after the MIS (detailed design
  % document) has been completed.}

% \wss{Reference your MIS (detailed design document) and explain your overall
% philosophy for test case selection.}  

% \wss{To save space and time, it may be an option to provide less detail in this section.  
% For the unit tests you can potentially layout your testing strategy here.  That is, you 
% can explain how tests will be selected for each module.  For instance, your test building 
% approach could be test cases for each access program, including one test for normal behaviour 
% and as many tests as needed for edge cases.  Rather than create the details of the input 
% and output here, you could point to the unit testing code.  For this to work, you code 
% needs to be well-documented, with meaningful names for all of the tests.}

% \subsection{Unit Testing Scope}

% \wss{What modules are outside of the scope.  If there are modules that are
%   developed by someone else, then you would say here if you aren't planning on
%   verifying them.  There may also be modules that are part of your software, but
%   have a lower priority for verification than others.  If this is the case,
%   explain your rationale for the ranking of module importance.}

% \subsection{Tests for Functional Requirements}

% \wss{Most of the verification will be through automated unit testing.  If
%   appropriate specific modules can be verified by a non-testing based
%   technique.  That can also be documented in this section.}

% \subsubsection{Module 1}

% \wss{Include a blurb here to explain why the subsections below cover the module.
%   References to the MIS would be good.  You will want tests from a black box
%   perspective and from a white box perspective.  Explain to the reader how the
%   tests were selected.}

% \subsubsection{Module 2}

% \subsection{Tests for Nonfunctional Requirements}

% \wss{If there is a module that needs to be independently assessed for
%   performance, those test cases can go here.  In some projects, planning for
%   nonfunctional tests of units will not be that relevant.}

% \wss{These tests may involve collecting performance data from previously
%   mentioned functional tests.}

% \subsection{Traceability Between Test Cases and Modules}

% \wss{Provide evidence that all modules have been considered.}



\newpage{}

\bibliographystyle{IEEEtran}
% \bibliographystyle{plainnat}
\bibliography{../../refs/References}

\section{Appendix}

This is where you can place additional information.

\subsection{Symbolic Parameters}

The following symbolic constants are used throughout this document:

\begin{table}[htbp]
	\centering
	\caption{Symbolic Constants Used in the Verification and Validation Plan}
	\label{tab:symbolicparameters}
	\begin{tabularx}{\textwidth}{| X | c | c |}
		\hline
		\textbf{Constant} & \textbf{Value} & \textbf{Unit}\\
		\hline
		\leavevmode\hypertarget{DISTINCT_LOCATION_ICONS}{DISTINCT\_LOCATION\_ICONS}
		 & 5 
		& Unitless/Quantity\\
		\hline
		\leavevmode\hypertarget{TRANSITION_TIME}{TRANSITION\_TIME} & 4 & 
		Seconds\\
		\hline
		\leavevmode\hypertarget{TRANSITION_TESTS}{TRANSITION\_TESTS} & 3 & 
		Unitless/Quantity\\
		\hline
	\end{tabularx}
\end{table}

\subsection{Usability Survey Questions?}

\wss{This is a section that would be appropriate for some projects.}

\newpage{}

\section{Appendix --- Generative AI Use Disclosure}

\subsection{AI Tools Used}

\begin{itemize}
	\item Anthropic Claude Sonnet 5.5
	\item Google Search AI Overview (Google Gemini)
\end{itemize}

\subsection{How and Where Used}

AI tools were used for assistance with LaTeX syntax and commands, during 
the development of sections SRS Verification and Design Verification, and when 
developing 
LaTeX formatting for the document as a whole. Specifically, AI tools were 
used for assistance in creating the LaTeX custom system test format macro 
used throughout the document, as well as the symbolic parameter table used in 
section 6.1. In these sections, 
AI tools were not used to generate any content text (LaTeX formatting code was 
generated), or wholly generate any of the 
ideas or conclusions presented.

\subsection{Verification of AI-Generated Content}

Any LaTeX syntax or code provided by AI tools was verified by testing 
compilation in LaTeX documents, and ensuring the resulting output was 
correctly formatted and matched the needs of the document. No ideas or text 
(outside of LaTeX formatting code) were 
newly generated by AI tools and presented in the document.

\newpage{}

\section{Appendix --- Reflection}

\textbf{Reflection Questions Answered:}
\begin{enumerate}
  \item What went well while writing this deliverable? 
  \item What pain points did you experience during this deliverable, and how
    did you resolve them?
  \item What knowledge and skills will the team collectively need to acquire to
  successfully complete the verification and validation of your project?
  Examples of possible knowledge and skills include dynamic testing knowledge,
  static testing knowledge, specific tool usage, Valgrind etc.  You should look to
  identify at least one item for each team member.
  \item For each of the knowledge areas and skills identified in the previous
  question, what are at least two approaches to acquiring the knowledge or
  mastering the skill?  Of the identified approaches, which will each team
  member pursue, and why did they make this choice?
\end{enumerate}

\textbf{Andreas:}
\begin{enumerate}
	\item 
	\item
	\item 
	\item 
\end{enumerate}

\textbf{Humna:}
\begin{enumerate}
	\item 
	\item
	\item 
	\item 
\end{enumerate}

\textbf{Christopher:}
\begin{enumerate}
	\item 
	\item
	\item 
	\item 
\end{enumerate}

\textbf{Owen:}
\begin{enumerate}
	\item 
	\item
	\item 
	\item 
\end{enumerate}

\textbf{Andy:}
\begin{enumerate}
	\item 
	\item
	\item 
	\item 
\end{enumerate}

\end{document}